% GENERATED FILE. Edit canonical lessons, then run npm run generate:books.
% canonical-source-hash: fnv1a64:3d99b4ed0c3edd11
% canonical-lessons: TE-C465-purapalaka-sangham, TE-C465-janabha, TE-C465-nivasi, TE-C465-sthanikudu, TE-C465-aparicitudu

\chapter{Townspeople}
\label{ch:te-townspeople}
\hlchaptermodality{465}

\begin{chapteropening}
\textbf{By the end of this chapter:} \emph{I can say a municipality, population, a resident, a local person and a stranger.}

Complete the last lesson of chapter 465: I can say a municipality, population, a resident, a local person and a stranger.
\end{chapteropening}

\section[\texorpdfstring{\te{పురపాలక} \te{సంఘం} (purapālaka saṅghaṁ)}{పురపాలక సంఘం (purapālaka saṅghaṁ)}]{\te{పురపాలక} \te{సంఘం} (purapālaka saṅghaṁ) — a municipality}

\label{lesson:TE-C465-purapalaka-sangham}



\begin{quote}
Before the new one: say the Telugu for a supply, then the Telugu for to collect payment.
\end{quote}

\subsection*{You'll want to know: \te{పురపాలక} \te{సంఘం}}
\textbf{\te{పురపాలక} \te{సంఘం}} — \emph{purapālaka saṅghaṁ} — \textquotedblleft{}a municipality\textquotedblright{}.

\begin{quote}
\te{పురపాలక} \te{సంఘం} \te{కొత్త} \te{రోడ్లు} \te{వేసింది}.
\end{quote}

\begin{quote}
\emph{purapālaka saṅghaṁ kotta rōḍlu vēsindi.}
\end{quote}

\begin{quote}
The municipality laid new roads.
\end{quote}

\te{పుర}, town, + \te{పాలక}, governing, + \te{సంఘం}, a body. A big city has a \te{నగరపాలక} \te{సంస్థ}.

\begin{cousinweb}
\textbf{\te{పురపాలక} \te{సంఘం}} (\emph{purapālaka saṅghaṁ}) is a learned Telugu phrase built from three Sanskrit words: \emph{pura}, \textquotedblleft{}a town, a fortress\textquotedblright{} (CDIAL 8278), \emph{pālaka}, \textquotedblleft{}guarding, governing\textquotedblright{} (Monier-Williams), and \emph{saṅgha}, \textquotedblleft{}a body, an association\textquotedblright{}.
\end{cousinweb}

\begin{grammarlens}[title={What you've built}]
One more word for getting around.
\end{grammarlens}

\subsection*{Your turn}
\begin{itemize}
\raggedright
  \item \emph{Say it:} \emph{purapālaka saṅghaṁ}
  \item \emph{Say it:} \emph{purapālaka saṅghaṁ}, once more
  \item \emph{Say it:} \emph{purapālaka saṅghaṁ kotta rōḍlu vēsindi.} — The municipality laid new roads.
  \item \emph{Recall:} say the Telugu for a supply, then the Telugu for to collect payment, then say \emph{purapālaka saṅghaṁ} again
\end{itemize}

\begin{tcolorbox}[breakable,colback=teal!4,colframe=teal!35!black,title={Before you move on}]
What does \te{పురపాలక} \te{సంఘం} mean? (\textquotedblleft{}A municipality\textquotedblright{}.) Say it once more.
\end{tcolorbox}
\section[\texorpdfstring{\te{జనాభా} (janābhā)}{జనాభా (janābhā)}]{\te{జనాభా} (janābhā) — population}

\label{lesson:TE-C465-janabha}



\begin{quote}
Before the new one: say the Telugu for to collect payment, then the Telugu for a municipality.
\end{quote}

\subsection*{You'll want to know: \te{జనాభా}}
\textbf{\te{జనాభా}} — \emph{janābhā} — \textquotedblleft{}population\textquotedblright{}.

\begin{quote}
\te{మా} \te{పట్టణం} \te{జనాభా} \te{లక్ష} \te{కంటే} \te{ఎక్కువ}.
\end{quote}

\begin{quote}
\emph{mā paṭṭaṇaṁ janābhā lakṣa kaṇṭē ekkuva.}
\end{quote}

\begin{quote}
The population of our town is more than a hundred thousand.
\end{quote}

\te{జనాభా} \te{లెక్కలు}, a census.

\begin{cousinweb}
\textbf{\te{జనాభా}} (\emph{janābhā}) is a Telugu word for population whose origin is not settled.
\end{cousinweb}

\begin{grammarlens}[title={What you've built}]
One more word for getting around.
\end{grammarlens}

\subsection*{Your turn}
\begin{itemize}
\raggedright
  \item \emph{Say it:} \emph{janābhā}
  \item \emph{Say it:} \emph{janābhā}, once more
  \item \emph{Say it:} \emph{mā paṭṭaṇaṁ janābhā lakṣa kaṇṭē ekkuva.} — The population of our town is more than a hundred thousand.
  \item \emph{Recall:} say the Telugu for to collect payment, then the Telugu for a municipality, then say \emph{janābhā} again
\end{itemize}

\begin{tcolorbox}[breakable,colback=teal!4,colframe=teal!35!black,title={Before you move on}]
What does \te{జనాభా} mean? (\textquotedblleft{}Population\textquotedblright{}.) Say it once more.
\end{tcolorbox}
\section[\texorpdfstring{\te{నివాసి} (nivāsi)}{నివాసి (nivāsi)}]{\te{నివాసి} (nivāsi) — a resident}

\label{lesson:TE-C465-nivasi}



\begin{quote}
Before the new one: say the Telugu for a municipality, then the Telugu for population.
\end{quote}

\subsection*{You'll want to know: \te{నివాసి}}
\textbf{\te{నివాసి}} — \emph{nivāsi} — \textquotedblleft{}a resident\textquotedblright{}.

\begin{quote}
\te{ఆ} \te{వీధి} \te{నివాసులు} \te{ఫిర్యాదు} \te{చేశారు}.
\end{quote}

\begin{quote}
\emph{ā vīdhi nivāsulu phiryādu cēśāru.}
\end{quote}

\begin{quote}
The residents of that street complained.
\end{quote}

Plural: \te{నివాసులు}, \emph{nivāsulu}.

Plural \te{నివాసులు}. \te{నివాసం} is a dwelling.

\begin{cousinweb}
\textbf{\te{నివాసి}} (\emph{nivāsi}) is borrowed from Sanskrit \emph{nivāsin}, \textquotedblleft{}dwelling, living in\textquotedblright{} (Monier-Williams), from \emph{ni-vas}, \textquotedblleft{}to dwell\textquotedblright{}.
\end{cousinweb}

\begin{grammarlens}[title={What you've built}]
One more word for getting around.
\end{grammarlens}

\subsection*{Your turn}
\begin{itemize}
\raggedright
  \item \emph{Say it:} \emph{nivāsi}
  \item \emph{Say it:} \emph{nivāsi}, once more
  \item \emph{Say it:} \emph{ā vīdhi nivāsulu phiryādu cēśāru.} — The residents of that street complained.
  \item \emph{Recall:} say the Telugu for a municipality, then the Telugu for population, then say \emph{nivāsi} again
\end{itemize}

\begin{tcolorbox}[breakable,colback=teal!4,colframe=teal!35!black,title={Before you move on}]
What does \te{నివాసి} mean? (\textquotedblleft{}A resident\textquotedblright{}.) Say it once more.
\end{tcolorbox}
\section[\texorpdfstring{\te{స్థానికుడు} (sthānikuḍu)}{స్థానికుడు (sthānikuḍu)}]{\te{స్థానికుడు} (sthānikuḍu) — a local person}

\label{lesson:TE-C465-sthanikudu}



\begin{quote}
Before the new one: say the Telugu for population, then the Telugu for a resident.
\end{quote}

\subsection*{You'll want to know: \te{స్థానికుడు}}
\textbf{\te{స్థానికుడు}} — \emph{sthānikuḍu} — \textquotedblleft{}a local person\textquotedblright{}.

\begin{quote}
\te{దారి} \te{తెలియక} \te{ఒక} \te{స్థానికుడిని} \te{అడిగాను}.
\end{quote}

\begin{quote}
\emph{dāri teliyaka oka sthānikuḍini aḍigānu.}
\end{quote}

\begin{quote}
Not knowing the way, I asked a local.
\end{quote}

Plural: \te{స్థానికులు}, \emph{sthānikulu}.

\te{స్థానం}, a place, + -\te{ఇకుడు}. \te{స్థానిక}, local, as an adjective: \te{స్థానిక} \te{వార్తలు}, local news.

\begin{cousinweb}
\textbf{\te{స్థానికుడు}} (\emph{sthānikuḍu}) is borrowed from Sanskrit \emph{sthānika}, \textquotedblleft{}belonging to a place, local\textquotedblright{} (Monier-Williams), from \emph{sthāna}, \textquotedblleft{}a place\textquotedblright{}.
\end{cousinweb}

\begin{grammarlens}[title={What you've built}]
One more word for getting around.
\end{grammarlens}

\subsection*{Your turn}
\begin{itemize}
\raggedright
  \item \emph{Say it:} \emph{sthānikuḍu}
  \item \emph{Say it:} \emph{sthānikuḍu}, once more
  \item \emph{Say it:} \emph{dāri teliyaka oka sthānikuḍini aḍigānu.} — Not knowing the way, I asked a local.
  \item \emph{Recall:} say the Telugu for population, then the Telugu for a resident, then say \emph{sthānikuḍu} again
\end{itemize}

\begin{tcolorbox}[breakable,colback=teal!4,colframe=teal!35!black,title={Before you move on}]
What does \te{స్థానికుడు} mean? (\textquotedblleft{}A local person\textquotedblright{}.) Say it once more.
\end{tcolorbox}
\section[\texorpdfstring{\te{అపరిచితుడు} (aparicituḍu)}{అపరిచితుడు (aparicituḍu)}]{\te{అపరిచితుడు} (aparicituḍu) — a stranger}

\label{lesson:TE-C465-aparicitudu}



\begin{quote}
Before the new one: say the Telugu for a resident, then the Telugu for a local person.
\end{quote}

\subsection*{You'll want to know: \te{అపరిచితుడు}}
\textbf{\te{అపరిచితుడు}} — \emph{aparicituḍu} — \textquotedblleft{}a stranger\textquotedblright{}.

\begin{quote}
\te{ఒక} \te{అపరిచితుడు} \te{నాకు} \te{దారి} \te{చూపించాడు}.
\end{quote}

\begin{quote}
\emph{oka aparicituḍu nāku dāri cūpiñcāḍu.}
\end{quote}

\begin{quote}
A stranger showed me the way.
\end{quote}

Plural: \te{అపరిచితులు}, \emph{aparicitulu}.

\te{అ}-, not, + \te{పరిచితుడు}, one who is known. \te{పరిచయం}, an acquaintance, is the same root.

\begin{cousinweb}
\textbf{\te{అపరిచితుడు}} (\emph{aparicituḍu}) is borrowed from Sanskrit \emph{aparicita}, \textquotedblleft{}unknown, unacquainted\textquotedblright{} (Monier-Williams), the negative of \emph{paricita}, \textquotedblleft{}known\textquotedblright{}.
\end{cousinweb}

\begin{grammarlens}[title={What you've built}]
One more word for getting around.
\end{grammarlens}

\subsection*{Your turn}
\begin{itemize}
\raggedright
  \item \emph{Say it:} \emph{aparicituḍu}
  \item \emph{Say it:} \emph{aparicituḍu}, once more
  \item \emph{Say it:} \emph{oka aparicituḍu nāku dāri cūpiñcāḍu.} — A stranger showed me the way.
  \item \emph{Recall:} say the Telugu for a resident, then the Telugu for a local person, then say \emph{aparicituḍu} again
\end{itemize}

\begin{tcolorbox}[breakable,colback=teal!4,colframe=teal!35!black,title={Before you move on}]
What does \te{అపరిచితుడు} mean? (\textquotedblleft{}A stranger\textquotedblright{}.) Say it once more.
\end{tcolorbox}
